%============================================================
\chapter{Fazit und Ausblick}
\label{chap:fazit}
%============================================================

\section{Zusammenfassung}
\label{sec:summary}

Im Rahmen des Projekts \emph{Mxck Follow the Gap} wurde der
Follow-the-Gap-Algorithmus auf einem F1TENTH-Fahrzeug implementiert
und erprobt.
Die theoretische Auseinandersetzung mit dem Verfahren sowie die
erfolgreichen Simulationstests bestätigen die grundsätzliche
Eignung des Ansatzes für reaktive Hindernissvermeidung.

Das übergeordnete Projektziel -- ein vollständig autonomes und
kollisionsfreies Durchfahren eines realen Hindernisparcours --
konnte jedoch nicht erreicht werden.
Die wesentlichen Ursachen lassen sich wie folgt bündeln:

\begin{enumerate}[label=\arabic*., leftmargin=*]
  \item \textbf{Sim-to-Real-Diskrepanz:}
        Parameter, die im Simulator gute Ergebnisse liefern, können
        auf der realen Plattform zu instabilem Verhalten führen.
  \item \textbf{Sensorqualität:}
        Bodenreflexionen des Linoleum\-bodens störten die
        LiDAR-Messungen und erzeugten Phantomhindernisse.
  \item \textbf{Parametertuning:}
        Der Sicherheitsblasen-Radius wurde für reale Bedingungen
        zu groß gewählt, was freie Lücken in engen Passagen blockierte.
  \item \textbf{Zeitrahmen:}
        Der vorhandene Projektzeitraum ließ keine ausreichend iterative
        Optimierung der Hardwareinstallation zu.
\end{enumerate}

%------------------------------------------------------------
\section{Empfehlungen für weiterführende Arbeiten}
\label{sec:future_work}
%------------------------------------------------------------

\begin{description}[leftmargin=*]
  \item[Systematisches Hardware-Tuning:]
        Parameter sollten direkt auf dem realen Fahrzeug in kontrollierten
        Umgebungen optimiert werden, z.\,B.\ mittels automatisierter
        Bayes-Optimierung \cite{snoek2012practical}.
  \item[Sensorvorverarbeitung:]
        Eine dedizierte Filterpipeline für Boden- und
        Decken\-reflexionen (z.\,B.\ Kombination aus IMU-Daten und
        LiDAR-Intensitätswerten) kann die Qualität der Messpunkte
        wesentlich verbessern.
  \item[Hybridansatz:]
        Die Kombination von FTG mit einer globalen Pfadplanung
        (z.\,B.\ $A^*$ oder RRT) könnte das reaktive Verfahren
        für komplexere Szenarien stabilisieren.
  \item[Genaueres Simulationsmodell:]
        Die Einbindung eines realistischeren Fahrzeugmodells
        (Reifenreibung, Motorträgheit, Latenzzeiten) in den
        Simulator würde die Übertragbarkeit der Parameter verbessern.
  \item[Tiefenkamera:]
        Die bereits verbaute Intel RealSense D435i könnte genutzt
        werden, um 3D-Informationen zu gewinnen und
        Boden\-reflexionen durch Höhenfilterung auszuschließen.
\end{description}

%------------------------------------------------------------
\section{Lernziele und Erkenntnisse}
\label{sec:lessons}
%------------------------------------------------------------

Trotz des nicht erreichten Projektziels konnten wertvolle
Erkenntnisse gewonnen werden:

\begin{itemize}
  \item Der \emph{Sim-to-Real-Gap} ist für reaktive Verfahren
        kritisch und muss beim Systemdesign frühzeitig
        berücksichtigt werden.
  \item Die Qualität der Sensordaten (Rauschen, Reflexionen,
        Auflösung) hat einen entscheidenden Einfluss auf die
        Robustheit des Algorithmus.
  \item Iteratives Testing und frühzeitiges Hardware-Debugging
        sind unerlässlich für den Projekterfolg.
  \item Die F1TENTH-Plattform bietet eine exzellente Grundlage für
        autonome Fahrexperimente und ist trotz der Herausforderungen
        für Lehrveranstaltungen gut geeignet.
\end{itemize}
